\section{刚体定轴转动}
\sourceproblem{选择1}{7}{13}
\problemtext 两均质圆盘质量和厚度相同，中心轴转动惯量 $J_B>J_A$，比较密度。
\workings $J=mR^2/2$，故 $R_B>R_A$；$\rho=m/(\pi R^2h)$，因而 $\rho_A>\rho_B$，选A。
\crosscheck 固定质量、厚度时 $J\rho=m^2/(2\pi h)$ 为常量。
\sourceproblem{选择2}{7}{13}
\problemtext 等质量、等半径的两个细环，一个质量分布均匀、另一个不均匀，比较中心垂直轴惯量。
\workings 每一质量元距轴均为 $R$，$J=\int R^2\,\mathrm dm=mR^2$，故 $J_1=J_2$，选C。
\crosscheck 结论针对细环几何中心的垂直轴，不依赖质心是否位于环心。
\sourceproblem{选择3}{7}{13}
\problemtext 按图1，圆盘顺时针转动，两等大反向平行力分别沿上、下两条作用线，判断角速度变化。
\workings 两力合力为零但构成力偶；设作用线垂直间距 $d>0$，图示合力矩为顺时针，大小 $Fd$。取顺时针为正，$\dot\omega=Fd/J>0$，故 $\omega_2>\omega_1$，选C。
\crosscheck 根据原图上方力向右、下方力向左判定力矩；不可用合力为零推出合力矩为零。
\sourceproblem{选择4}{7}{13}
\problemtext 均质杆长 $L$、质量 $M$，从水平静止绕端点水平轴下摆到竖直，判断力矩、角速度、角加速度变化。
\workings 设杆下摆角 $\phi$ 从水平起算，$\tau=MgL\cos\phi/2$，$\alpha=3g\cos\phi/(2L)$，随 $\phi$ 从0到 $\pi/2$ 均减小；能量守恒 $\omega^2=3g\sin\phi/L$，角速率增大，选A。
\crosscheck 竖直最低点角加速度为零但角速率最大。
\sourceproblem{选择5}{7}{13}
\problemtext 相同滑轮A悬挂质量 $m$，B直接受恒力 $G=mg$，比较角加速度。
\workings 轮半径 $R$、惯量 $J>0$；A有 $mg-T=ma$，$TR=J a/R$，故 $\beta_A=mgR/(J+mR^2)$；B有 $\beta_B=mgR/J$，故 $\beta_A<\beta_B$，选C。
\crosscheck A的张力小于 $mg$，因此不能把悬挂重物等同于恒拉力。
\sourceproblem{选择6}{8}{14}
\problemtext 图4中左侧 $m_1$ 小于右侧 $m_2$，滑轮某刻逆时针转动，比较绳两侧张力。
\workings 取右侧向下为正，$a=(m_2-m_1)g/(m_1+m_2+J/R^2)>0$。轮的角加速度为顺时针，$T_2-T_1=Ja/R^2>0$，故右侧张力较大，选C。
\crosscheck 当前逆时针角速度与顺时针角加速度并不矛盾；题设有惯量的轮两侧张力不等。
\sourceproblem{选择7}{8}{14}
\problemtext 固定光滑轴刚体受到数个合力为零的力，判断转速是否变化。
\workings 合力为零不决定轴向合力矩；$\dot\omega=\tau_z/J$。若 $\tau_z=0$ 则转速不变，否则通常改变，选D。
\crosscheck 力偶提供反例：净力为零但力矩非零。
\sourceproblem{选择8}{8}{14}
\problemtext 图5中两同质量、同速率、相向且共线的物体嵌入绕固定轴转动的圆盘，比较碰撞前后角速率。
\workings 两入射物体的轴向角动量互相抵消。短时碰撞中忽略外轴向冲量矩，$J\omega_1=(J+\Delta J)\omega_2$，其中 $\Delta J>0$，故 $\omega_2<\omega_1$（初始角速率非零），选A。
\crosscheck 不应用机械能守恒；嵌入是非弹性碰撞。
\sourceproblem{选择9}{8}{14}
\problemtext 比较杆与小球的非弹性碰撞系统A、人及自由旋转盘系统B的守恒量。
\workings 碰撞短时极限中A对轴的外冲量矩为零；B在忽略轴摩擦时外轴向力矩为零，故两者均有轴向角动量守恒。A的碰撞会耗散机械能；B可由人体内能转化为机械能。固定轴支承可给外冲量，线动量并不普遍守恒，选B。
\crosscheck A的角动量守恒只针对碰撞短时过程，非其后摆动全过程。
\sourceproblem{选择10}{8}{14}
\problemtext 自由绕固定光滑轴转动的物体热膨胀，判断角速度。
\workings 无外轴向力矩时 $J\omega$ 守恒。通常均匀各向同性热膨胀使各质量元轴距按比例增大，$J$ 增大、角速率减小，选B。
\crosscheck 若只沿轴向膨胀或材料各向异性，未必增加 $J$；原选项B采用常规径向膨胀假设。
